\documentclass[11pt]{article}
\usepackage[a4paper,margin=26mm]{geometry}
\usepackage[T1]{fontenc}
\usepackage{lmodern}
\usepackage{amsmath,amssymb,amsthm}
\usepackage{microtype}
\usepackage{array}
\usepackage{xurl}
\usepackage{hyperref}
\hypersetup{colorlinks=true,linkcolor=blue,urlcolor=blue,citecolor=blue,
 pdftitle={The largest domination roots of friendship graphs and of joins with repeated petals},
 pdfauthor={Horace Dong}}
\newtheorem{theorem}{Theorem}
\newtheorem{lemma}{Lemma}
\newtheorem{proposition}{Proposition}
\newtheorem{corollary}{Corollary}
\theoremstyle{definition}
\newtheorem{definition}{Definition}
\theoremstyle{remark}
\newtheorem{remark}{Remark}
\newcommand{\repoversion}{1.25.0}
\setlength{\emergencystretch}{1.5em}
\newcommand{\C}{\mathbb C}
\newcommand{\Log}{\operatorname{Log}}

\title{The largest domination roots of friendship graphs\\ and of joins with repeated petals}
\author{Horace Dong\\ \small Independent researcher\\
\small\texttt{maxwelldhx@gmail.com}\\
\small ORCID: \href{https://orcid.org/0009-0002-7554-7548}{0009-0002-7554-7548}}
\date{22 November 2026}
\begin{document}
\maketitle

\begin{abstract}
The domination polynomial $D(G,x)$ of a graph $G$ counts the dominating sets
of $G$ by size. Alikhani, Brown and Jahari asked for a good upper bound on the
modulus of the domination roots of the friendship graph $F_m=K_1\vee mK_2$.
Fix a simple graph $B$ with at least one vertex and an integer $a\ge1$, and let $J_m$ be the
join of the complete graph $K_a$ with $m$ disjoint copies of $B$. We show that
the largest modulus $M_m$ of a complex root of $D(J_m,x)$ is attained and
satisfies $M_m\sim\bigl(c\,s\,m/(a\log m)\bigr)^{1/c}$ as $m\to\infty$, where
$c=1+\delta(B)$ and $s$ is the number of distinct closed neighbourhoods of size
$c$ in $B$. For friendship graphs this gives $M_m\sim\sqrt{2m/\log m}$. In
this case the upper half of the asymptotics already follows, by a short
computation, from a recent implicit bound of Rather; what is added there is
the matching lower bound, which shows that Rather's implicit bound is
asymptotically sharp, while his explicit bound is not. The convergence is
slow. The three main statements are formally verified in Lean~4. The proofs
were found and formalized by AI agents working under the author's direction;
their role is described at the end of the paper.
\end{abstract}

\section{Introduction}

All graphs are finite, simple and undirected. A set $S$ of vertices of a graph
$G$ is \emph{dominating} if every vertex of $G$ lies in $S$ or has a neighbour
in $S$. The \emph{domination polynomial} of $G$ is
\[
D(G,x)=\sum_{S\ \text{dominating}} x^{|S|}=\sum_{i\ge0} d(G,i)\,x^i ,
\]
where $d(G,i)$ is the number of dominating sets of size $i$; its roots are the
\emph{domination roots} of $G$. The \emph{friendship graph} $F_m$ consists of
$m$ triangles with one common vertex; equivalently $F_m=K_1\vee mK_2$, the
join of one vertex with $m$ disjoint edges. It has $2m+1$ vertices, and
\begin{equation}\label{eq:friend}
D(F_m,x)=(x^2+2x)^m+x(1+x)^{2m}
\end{equation}
\cite[Theorem~2.1]{abj}. Alikhani, Brown and Jahari showed that the limits of
the domination roots of $F_m$ are $0$ and a hyperbola~\cite[Theorem~2.5]{abj},
so that the moduli of the roots are unbounded, and asked:

\begin{quote}
\textbf{Question 3.2 of~\cite{abj}.} What is a good upper bound on the
modulus of the roots of $F_n$?
\end{quote}

\noindent They added that some calculations seemed to indicate that the
moduli grow quite slowly, ``perhaps like ln n''~\cite[p.~176]{abj}. We use
the numbering of the journal version~\cite{abj}, where the question is on
p.~176. In the arXiv version (arXiv:1401.2092v2) the same question is
Question~2 of Section~3, with the same wording, and Theorems~2.1 and~2.5 are
Theorems~1 and~4 (the statement of the latter contains misprints).

Rather~\cite[Section~4]{ra} proved that every domination root $z$ of $F_m$
with $|z|>1$ satisfies $(|z|-1)^2\log|z|\le m$ \cite[Theorem~4.1]{ra}. He
deduced the explicit bound $|z|\le 1+\sqrt{m/\log 2}$
\cite[Corollary~4.2]{ra} and noted the sharper implicit bound $|z|\le R_m$,
where $R_m>1$ is the solution of $(R_m-1)^2\log R_m=m$ \cite[Remark~4.3]{ra}.
He wrote that his estimate is ``not expected to be asymptotically
sharp''~\cite[p.~10]{ra}, and in his conclusion~\cite[Section~7]{ra} he left open whether
the
modulus bound for friendship graphs is asymptotically optimal.

We determine the first-order asymptotics of the largest modulus, for the
friendship graphs and for a larger family. Let $B$ be a fixed graph with at
least one vertex (it may be disconnected and may have isolated vertices), and
let $a\ge1$ be a fixed integer. For $m\ge1$ let
\[
J_m=J_m(a,B)=K_a\vee mB
\]
be the graph obtained from a complete graph $K_a$ (the \emph{centre}) and $m$
disjoint copies of $B$ (the \emph{petals}) by joining every centre vertex to
every petal vertex. Let $b=|V(B)|$, let $\delta(B)$ be the minimum degree of
$B$, and put
\[
c=1+\delta(B),\qquad s=\#\{N_B[v]:v\in V(B),\ |N_B[v]|=c\},
\]
where $N_B[v]$ is the closed neighbourhood of $v$ in $B$; thus $c$ is the
smallest size of a closed neighbourhood, and $s$ counts the \emph{distinct}
closed neighbourhoods of that size (a set that is the closed neighbourhood of
several vertices is counted once). For $m\ge2$ let
\[
S_m=\Bigl(\frac{c\,s\,m}{a\log m}\Bigr)^{1/c},
\]
with the natural logarithm. The friendship graph is $F_m=J_m(1,K_2)$; for
$B=K_2$ both closed neighbourhoods are the whole vertex set, so $c=2$, $s=1$
and $S_m=\sqrt{2m/\log m}$.

\begin{theorem}\label{thm:attain}
For every graph $B$ with at least one vertex and all integers $a,m\ge1$, the
set $\{|z|: z\in\C,\ D(J_m,z)=0\}$ has a largest element. We denote it by
$M_m=M_m(a,B)$.
\end{theorem}

\begin{theorem}\label{thm:main}
For every graph $B$ with at least one vertex and every integer $a\ge1$,
\[
\lim_{m\to\infty}\frac{M_m(a,B)}{S_m}=1 .
\]
Equivalently, for every $\varepsilon>0$ there is $m_0$ such that for every
$m\ge m_0$ every domination root $z$ of $J_m$ satisfies
$|z|\le(1+\varepsilon)S_m$, and some domination root $z$ of $J_m$ satisfies
$|z|\ge(1-\varepsilon)S_m$.
\end{theorem}

\begin{corollary}\label{cor:friend}
The largest modulus $M_m$ of a domination root of the friendship graph $F_m$
satisfies $M_m/\sqrt{2m/\log m}\to1$ as $m\to\infty$.
\end{corollary}

\noindent Theorems~\ref{thm:attain} and~\ref{thm:main} (in the limit form)
and Corollary~\ref{cor:friend} are the three statements verified in Lean
(Section~\ref{sec:lean}); the $\varepsilon$-form of Theorem~\ref{thm:main} is
equivalent to the limit form by Theorem~\ref{thm:attain} and is what the Lean
proof establishes on the way. Remark~\ref{rem:rather}, the examples of
$(c,s)$ other than $B=K_2$, the consequence $M_m/\log m\to\infty$ and the
numerical values in Table~\ref{tab:num} are not part of the Lean statements.

Thus Question~3.2 of~\cite{abj} has a sharp answer in the asymptotic sense:
the roots of $F_m$ have modulus at most $(1+o(1))\sqrt{2m/\log m}$, and no
upper bound of the form $(1-\varepsilon)\sqrt{2m/\log m}$ holds for large $m$.
This does not give an explicit bound for each fixed $m$; for those, Rather's
bounds remain the available ones. By Corollary~\ref{cor:friend},
$M_m/\log m\to\infty$, so the
moduli grow faster than the logarithmic growth suggested by the numerical
data in~\cite{abj} (a consequence of Corollary~\ref{cor:friend} that is not
stated separately in Lean).

Some further examples of $(c,s)$; Theorem~\ref{thm:main} applies to each of
them, and the values of $(c,s)$ were computed by hand (only for $B=K_2$ is
this evaluation also done in Lean): $B=K_1$ gives $c=s=1$ and $S_m=m/(a\log m)$; the graph of $b$
isolated vertices gives $c=1$, $s=b$; $B=K_b$ gives $c=b$, $s=1$, so for the
windmill graphs $K_1\vee mK_b$ the largest modulus is asymptotic to
$(bm/\log m)^{1/b}$; the path $P_3$ gives $(2,2)$, the star $K_{1,3}$ gives
$(2,3)$ and the cycle $C_4$ gives $(3,4)$. For $B=K_3$, counting vertices of
minimum degree instead of distinct closed neighbourhoods would give $s=3$
instead of $s=1$ and the wrong constant.

\subsection*{Prior work and what this paper adds}
The friendship formula~\eqref{eq:friend}, the limiting hyperbola and the
question are from~\cite{abj}. The formula for $D(J_m,x)$ used below
(Lemma~\ref{lem:join}) follows from standard facts: the domination polynomial
is multiplicative over disjoint unions~\cite[Theorem~1]{ap},
\cite[Section~1]{kpstt}, and the join satisfies
$D(G_1\vee G_2,x)=\bigl((1+x)^{n_1}-1\bigr)\bigl((1+x)^{n_2}-1\bigr)+D(G_1,x)+D(G_2,x)$
for graphs of orders $n_1$ and $n_2$~\cite[Theorem~2]{ap}; both facts are
used in the proof of~\cite[Theorem~2.1]{abj}, and for $a=1$ the formula also
follows from the corona formula~\cite[Theorem~3.12]{kpstt}, since
$K_1\vee mB$ is the corona $K_1\circ mB$. Only the
Lean derivation of Lemma~\ref{lem:join} from the definition is ours.

For friendship graphs, Rather's implicit bound $R_m$ already gives the upper
half of Corollary~\ref{cor:friend}, because $R_m\sim\sqrt{2m/\log m}$
(Remark~\ref{rem:rather}; this short computation is not in~\cite{ra} and not
in Lean). The parts not covered by~\cite{ra} or by the other work we found are: the
matching lower bound, that is, actual roots
of modulus $(1-o(1))\sqrt{2m/\log m}$ for every large $m$, which shows that
$M_m/R_m\to1$; and Theorem~\ref{thm:main} for every fixed petal $B$ and every
fixed centre $K_a$, with both bounds. A phenomenon of the same kind is known
for chromatic polynomials: Brown, Hickman, Sokal and Wagner~\cite{bhsw}
proved that the chromatic roots of $k$-ary generalized theta graphs lie in
$|z-1|\le(1+o(1))\,k/\log k$ and that $K_{2,k}$ has a chromatic root of that
modulus. Searches up to 20 November 2026 found no earlier sharp asymptotic
for the largest domination roots of friendship graphs and no matching lower
bound (see the final section).

\section{From dominating sets to a root equation}\label{sec:equation}

Throughout, $B$, $a$, $b$, $c$, $s$ are as in the introduction, and all
polynomials have complex coefficients.

\begin{lemma}[join formula]\label{lem:join}
For $m\ge1$,
\[
D(J_m,z)=\bigl((1+z)^a-1\bigr)(1+z)^{bm}+D(B,z)^m .
\]
\end{lemma}

\begin{proof}
Let $S$ be a set of vertices of $J_m$. If $S$ contains a centre vertex, then
$S$ is dominating, since a centre vertex is adjacent to every other vertex;
these sets contribute $((1+z)^a-1)(1+z)^{bm}$. If $S$ contains no centre
vertex, then a petal vertex is dominated only from its own copy of $B$, so
$S$ is dominating if and only if its intersection with each copy is a
dominating set of that copy. In this case $S$ is nonempty (a dominating set
of a graph with a vertex is nonempty, and $m\ge1$), so it also dominates the
centre. Choosing the intersections independently for the $m$ copies gives the
contribution $D(B,z)^m$.
\end{proof}

\paragraph{Bad sets.}
Call a set $T\subseteq V(B)$ \emph{bad} if $N_B[v]\subseteq T$ for some vertex
$v$. A set $S$ fails to dominate $B$ exactly when some vertex has no element
of $N_B[v]$ in $S$, that is, when $V(B)\setminus S$ is bad. Every bad set has
at least $c$ elements, and the bad sets with exactly $c$ elements are exactly
the $s$ distinct closed neighbourhoods of size $c$. Define
\[
Q(w)=1-\sum_{T\ \text{bad}} w^{|T|}(1-w)^{b-|T|}.
\]

\begin{lemma}[leading missing term]\label{lem:Q}
\begin{enumerate}
\item[(a)] $Q(w)=\sum_{S\ \text{dominating in}\ B}(1-w)^{|S|}w^{b-|S|}$.
\item[(b)] For $y\ne0$, $y^b\,Q(1/y)=D(B,y-1)$.
\item[(c)] There is a polynomial $E$ such that
$Q(w)=1-s\,w^c+w^{c+1}E(w)$.
\end{enumerate}
\end{lemma}

\begin{proof}
(a) The weights $(1-w)^{|S|}w^{b-|S|}$ of all subsets $S$ of $V(B)$ add up to
$((1-w)+w)^b=1$. The non-dominating sets are the complements of bad sets, and
the complement of $T$ has weight $w^{|T|}(1-w)^{b-|T|}$; subtracting their
weights from $1$ leaves the sum over dominating sets. (b) By (a),
$y^bQ(1/y)=\sum_S (1-1/y)^{|S|}y^{-(b-|S|)}y^b=\sum_S (y-1)^{|S|}$, the sum over
dominating sets $S$. (c) A bad set $T$ with $|T|=c$ contributes
$w^c(1-w)^{b-c}=w^c-w^{c+1}(\cdots)$, and there are $s$ of them; a bad set with
$|T|>c$ contributes a multiple of $w^{c+1}$.
\end{proof}

\begin{lemma}[root equation]\label{lem:rooteq}
Let $m\ge1$ and $y\ne0$. Then $y-1$ is a domination root of $J_m$ if and only
if
\[
Q(1/y)^m=1-y^a .
\]
\end{lemma}

\begin{proof}
By Lemmas~\ref{lem:join} and~\ref{lem:Q}(b),
$D(J_m,y-1)=(y^a-1)y^{bm}+\bigl(y^bQ(1/y)\bigr)^m
=y^{bm}\bigl(y^a-1+Q(1/y)^m\bigr)$, and $y^{bm}\ne0$.
\end{proof}

\begin{proof}[Proof of Theorem~\ref{thm:attain}]
For any graph $G$ with $n\ge1$ vertices, the only dominating set of size $n$
is the whole vertex set, so $D(G,x)$ has leading coefficient $1$ in degree
$n$ and is a nonzero polynomial; its root set is therefore finite. The empty
set dominates no graph with a vertex, so $D(G,0)=0$ and the root set is
nonempty. A finite nonempty set of moduli has a largest element. Apply this
to $G=J_m$, which has $a+bm\ge1$ vertices.
\end{proof}

\section{The upper bound}\label{sec:upper}

\begin{lemma}[the scale]\label{lem:scale}
For $m\ge2$, $S_m>0$ and $S_m^c=c\,s\,m/(a\log m)$. Moreover
$\log S_m/\log m\to1/c$ and $S_m\to\infty$ as $m\to\infty$.
\end{lemma}

\begin{proof}
The first two claims are immediate. Since
$c\log S_m=\log(cs/a)+\log m-\log\log m$ and $\log\log m/\log m\to0$, we get
$\log S_m/\log m\to1/c$, and then $\log S_m\to\infty$.
\end{proof}

\begin{proposition}[upper bound]\label{prop:upper}
For every $\varepsilon>0$ there is $m_0$ such that for every $m\ge m_0$ every
domination root $z$ of $J_m$ satisfies $|z|\le(1+\varepsilon)S_m$.
\end{proposition}

\begin{proof}
Let $E$ be as in Lemma~\ref{lem:Q}(c), let $K\ge0$ bound $|E(w)|$ on the
closed unit disc, and let $\beta=1+\varepsilon/2$ and $T=T_m=\beta S_m$.
Consider $y\in\C$ with $|y|=r\ge T$, where $m$ is so large that $T\ge1$ and
$T^a\ge2$. With
$w=1/y$ we have $|w|\le1/T\le1$, so
\[
|Q(1/y)|\le 1+|w|^c\bigl(s+|w|\,|E(w)|\bigr)\le 1+\frac{s+K/T}{T^c},
\qquad
|Q(1/y)|^m\le\exp\Bigl(\frac{m\,(s+K/T)}{T^c}\Bigr).
\]
By Lemma~\ref{lem:scale}, $m\,s/T^c=a\log m/(c\beta^c)$ and
$m\,K/T^{c+1}=(K/(sT))\,a\log m/(c\beta^c)$, so the exponent divided by
$\log m$ tends to $a/(c\beta^c)$. On the other hand $\log(T^a/2)/\log m\to a/c$,
which is strictly larger because $\beta>1$. Hence for all large $m$,
\[
|Q(1/y)|^m<T^a/2\le T^a-1\le r^a-1\le|1-y^a| ,
\]
so $y$ does not satisfy the root equation of Lemma~\ref{lem:rooteq}. Therefore
every domination root $z=y-1$ of $J_m$ with $y\ne0$ has $|y|<T$, and the root
$z=-1$ (if it is one) has $|z|=1$. In both cases
$|z|\le|y|+1\le\beta S_m+1\le(1+\varepsilon)S_m$ as soon as
$S_m\ge2/\varepsilon$, which holds for large $m$ by Lemma~\ref{lem:scale}.
\end{proof}

\section{Roots of the predicted size}\label{sec:lower}

\begin{proposition}[lower bound]\label{prop:lower}
For every $\varepsilon>0$ there is $m_0$ such that for every $m\ge m_0$ some
domination root $z$ of $J_m$ satisfies $|z|\ge(1-\varepsilon)S_m$.
\end{proposition}

The idea is to look for roots $y=S_m\omega v$ of the root equation, where
$\omega^c=-1$ and $v$ is close to~$1$. After taking logarithms and dividing by
$\log m$, the root equation becomes $H_m(v)=0$ for analytic functions $H_m$
that converge to $H(v)=\frac ac(v^{-c}-1)$, which vanishes at $v=1$; a
maximum-modulus argument then produces an actual zero of $H_m$ near~$1$.

\paragraph{The functions $H_m$.}
Fix $\omega\in\C$ with $\omega^c=-1$ (so $|\omega|=1$), and let $\Log u=\log|u|+i\operatorname{Arg}u$ with
$\operatorname{Arg}u\in(-\pi,\pi]$ be the principal logarithm of $u\ne0$ (as in
Mathlib); it satisfies $\exp(\Log u)=u$ and is holomorphic on the slit plane
$\C\setminus(-\infty,0]$. Write $Q(w)=1+\Delta(w)$ with
$\Delta(w)=w^c\bigl(-s+wE(w)\bigr)$, as in Lemma~\ref{lem:Q}(c). For $m\ge2$
and $v\ne0$ put $w=w_m(v)=(S_m\omega v)^{-1}$ and
\[
H_m(v)=\frac{m}{\log m}\Log\bigl(1+\Delta(w)\bigr)-a\,\frac{\log S_m}{\log m}
-\frac{a\Log v+\Log(-\omega^a)+\Log(1-w^a)}{\log m},
\]
for $v\ne0$; it is holomorphic in $v$ wherever $v$, $1+\Delta(w)$ and $1-w^a$
lie in the slit plane (the constant $\Log(-\omega^a)$ needs no such
condition).

\begin{lemma}[recovering a root]\label{lem:recover}
If $v$, $1+\Delta(w)$ and $1-w^a$ lie in the slit plane and $H_m(v)=0$,
then $y=S_m\omega v$ satisfies $Q(1/y)^m=1-y^a$.
\end{lemma}

\begin{proof}
Multiply $H_m(v)=0$ by $\log m$ and exponentiate: since $\exp(m\Log u)=u^m$ and
$\exp(a\Log v)=v^a$,
$Q(w)^m=S_m^a\,v^a\,(-\omega^a)(1-w^a)=-y^a(1-y^{-a})=1-y^a$, and $w=1/y$.
\end{proof}

\begin{lemma}[joint limit]\label{lem:limit}
Let $p\ne0$ lie in the slit plane. As $m\to\infty$ and $v\to p$ jointly,
$H_m(v)\to\frac ac\,(p^{-c}-1)$.
\end{lemma}

\begin{proof}
As $m\to\infty$ and $v\to p$, $w\to0$ and hence $\Delta(w)\to0$. Write
$\Log(1+t)=t\,L(t)$ with $L(t)\to1$ as $t\to0$. By Lemma~\ref{lem:scale},
$\frac m{\log m}w^c=\frac m{\log m}(\omega v)^{-c}S_m^{-c}
=\frac a{cs}(\omega v)^{-c}$, so the first term equals
$\frac a{cs}(\omega v)^{-c}\bigl(-s+wE(w)\bigr)L(\Delta(w))
\to-\frac ac\,\omega^{-c}p^{-c}=\frac ac\,p^{-c}$. The second term tends to
$-a/c$ by Lemma~\ref{lem:scale}. In the third term the numerator tends to the
finite value $a\Log p+\Log(-\omega^a)$, so the term tends to~$0$.
\end{proof}

\begin{lemma}[a zero from a dip]\label{lem:dip}
Let $F_m$ be functions that are holomorphic on an open disc
$U=\{|v-v_0|<r\}$ and continuous on its closure, for all large $m$. Suppose
$F_m(v_0)\to0$ and, for every $u$ with $|u-v_0|=r$, $F_m(v)$ tends to a
nonzero limit as $m\to\infty$ and $v\to u$ jointly. Then for all large $m$,
$F_m$ has a zero in $\overline U$.
\end{lemma}

\begin{proof}
By the joint limits and compactness of the circle (finitely many
neighbourhoods of points of the circle suffice), there is $\eta>0$ such that
for all large $m$, $|F_m(u)|\ge\eta$ on the circle while $|F_m(v_0)|<\eta$;
thus $|F_m(v_0)|<|F_m(u)|$ for every $u$ on the circle. If such an $F_m$ had no
zero in $\overline U$, then $1/F_m$ would be holomorphic on $U$, continuous on
$\overline U$, and larger in modulus at the centre than everywhere on the
boundary, contradicting the maximum modulus principle.
\end{proof}

\begin{proof}[Proof of Proposition~\ref{prop:lower}]
Choose $r$ with $0<r<\varepsilon/2$ such that the closed disc
$\{|v-1|\le r\}$ lies in the slit plane and the circle $|v-1|=r$ contains no
$c$-th root of unity (there are finitely many, and $1$ is the centre). Since
$w_m(v)\to0$ uniformly on the disc, for all large $m$ the numbers
$1+\Delta(w)$ and $1-w^a$ lie in the slit plane for every $v$ in the closed
disc, so $H_m$ is holomorphic on a neighbourhood of it. By
Lemma~\ref{lem:limit}, $H_m(1)\to0$, and at each point $u$ of the circle
$H_m(v)\to\frac ac(u^{-c}-1)\ne0$ jointly. By Lemma~\ref{lem:dip}, for all
large $m$ there is $v$ with $|v-1|\le r$ and $H_m(v)=0$. By
Lemmas~\ref{lem:recover} and~\ref{lem:rooteq}, $z=S_m\omega v-1$ is a
domination root of $J_m$, and
$|z|\ge S_m|v|-1\ge(1-\varepsilon/2)S_m-1\ge(1-\varepsilon)S_m$ once
$S_m\ge2/\varepsilon$.
\end{proof}

\begin{proof}[Proof of Theorem~\ref{thm:main} and Corollary~\ref{cor:friend}]
By Theorem~\ref{thm:attain}, $M_m$ is the modulus of an actual root, so
Propositions~\ref{prop:upper} and~\ref{prop:lower} give
$1-\varepsilon\le M_m/S_m\le1+\varepsilon$ for all large $m$. For the
corollary take $a=1$ and $B=K_2$, where $c=2$, $s=1$ and
$S_m=\sqrt{2m/\log m}$.
\end{proof}

\section{Comparison with Rather's bounds, and numerics}\label{sec:rather}

\begin{remark}[not in Lean]\label{rem:rather}
Let $f(t)=(t-1)^2\log t$, which is strictly increasing on $(1,\infty)$, let
$R_m$ be defined by $f(R_m)=m$ as in~\cite[Remark~4.3]{ra}, and let
$T_m=\sqrt{2m/\log m}$. For fixed $\lambda>0$ we have
$(\lambda T_m-1)^2=\lambda^2\,\frac{2m}{\log m}(1+o(1))$ and
$\log(\lambda T_m)=\frac12\log m\,(1+o(1))$, so $f(\lambda T_m)/m\to\lambda^2$.
Taking $\lambda=1\pm\varepsilon$ with $0<\varepsilon<1$ and using monotonicity
gives
$(1-\varepsilon)T_m<R_m<(1+\varepsilon)T_m$ for large $m$, so
$R_m\sim T_m$. Together with Corollary~\ref{cor:friend}:
\begin{enumerate}
\item[(i)] Rather's implicit bound $|z|\le R_m$ already implies
$\limsup M_m/T_m\le1$, the upper half of Corollary~\ref{cor:friend};
\item[(ii)] this implicit bound is asymptotically sharp: $M_m/R_m\to1$;
\item[(iii)] the explicit bound of~\cite[Corollary~4.2]{ra} is not:
$(1+\sqrt{m/\log2})/M_m\sim\sqrt{\log m/(2\log2)}\to\infty$.
\end{enumerate}
So the question in~\cite{ra} whether the modulus bound is asymptotically
optimal has the answer ``yes'' for the implicit bound $R_m$ of Remark~4.3 and
``no'' for the explicit bound of Corollary~4.2.
\end{remark}

\paragraph{Slow convergence.}
The convergence in Corollary~\ref{cor:friend} is slow, so tables for small
$m$ can be misleading about the constant. Table~\ref{tab:num} lists, for the
friendship graphs, the ratios of $M_m$ and of $R_m$ to $T_m=\sqrt{2m/\log m}$.
For $m\le80$, $M_m$ was computed from all roots of~\eqref{eq:friend}; the
values $M_{10}=3.0650$ and $M_{20}=3.8428$ agree with Table~2 of~\cite{ra}.
For $m\ge10^4$ the entry is the modulus of a root of the root equation found
by Newton's method near the predicted location, which is a lower bound for
$M_m$; at $m=80$ the same method finds the root of largest modulus. At
every tabulated $m$ from $20$ to $10^8$ the computed ratio lies between $1.05$
and $1.06$; it approaches $1$ only slowly; the values are consistent with a relative error of order
$\log\log m/\log m$, which we have not proved. All these values were
computed in floating-point arithmetic (the final section says by whom) and
are not part of the Lean development.

\begin{table}[ht]
\centering\small
\begin{tabular}{@{}lcccccccccc@{}}
\hline
$m$ & $5$ & $10$ & $20$ & $40$ & $80$ & $10^4$ & $10^8$ & $10^{20}$ &
$10^{50}$ & $10^{100}$\\
\hline
$M_m/T_m$ & $1.004$ & $1.040$ & $1.052$ & $1.055$ & $1.055$ & $1.059$ &
$1.054$ & $1.034$ & $1.018$ & $1.010$\\
$R_m/T_m$ & $1.244$ & $1.273$ & $1.263$ & $1.240$ & $1.214$ & $1.103$ &
$1.063$ & $1.035$ & $1.018$ & $1.010$\\
\hline
\end{tabular}
\caption{Friendship graphs: the largest modulus $M_m$ and Rather's implicit
bound $R_m$, divided by $T_m=\sqrt{2m/\log m}$ (floating point; for
$m\ge10^4$ the first row uses the root found by Newton's method). Not
verified in Lean.}
\label{tab:num}
\end{table}

{\sloppy
\section{Formal verification}\label{sec:lean}

Theorems~\ref{thm:attain} and~\ref{thm:main} and Corollary~\ref{cor:friend}
are formally verified in Lean~4 (version 4.34.1) with Mathlib~\cite{mathlib}
(commit \texttt{d13f23b7}). The $14$ files are in version~\repoversion{} of
the repository~\cite{repo}, directory
\texttt{lean/Research/TaskP/RootRadius/}.

\paragraph{The statements.}
The statement file \texttt{Challenge.lean} contains only definitions and
statements, no proofs apart from the two short proofs that the adjacency
relation of \texttt{petalJoin} is symmetric and irreflexive. In ASCII
transliteration (the \texttt{import}, \texttt{open}, \texttt{namespace} and
\texttt{end} lines, the module docstring and the other docstrings omitted;
the file itself uses Unicode, e.g.\ $\mathbb C$, $\mathbb N$, $\mathbb R$, $\sum$, $\forall$,
$\exists$, $\in$, $\to$, $\vee$, $\wedge$, $\ne$, $\mathcal N$, $\oplus$, $\times$, $\langle\,\rangle$, $\top$ and
$\|z\|$):
{\footnotesize
\begin{verbatim}
def Dominates {V : Type*} (G : SimpleGraph V) (S : Finset V) : Prop :=
  forall v, v in S \/ exists u in S, G.Adj v u
noncomputable def dominationPolynomial {V : Type*} [Fintype V]
    (G : SimpleGraph V) : Complex[X] := by
  classical
  exact sum S : Finset V, if Dominates G S then X ^ S.card else 0
def petalJoin {V : Type*} (a m : Nat) (B : SimpleGraph V) :
    SimpleGraph (Fin a (+) (Fin m x V)) where
  Adj
    | .inl i, .inl j => i != j
    | .inl _, .inr _ => True
    | .inr _, .inl _ => True
    | .inr x, .inr y => x.1 = y.1 /\ B.Adj x.2 y.2
  symm := by
    constructor
    rintro (i | x) (j | y) h
    . exact Ne.symm h
    . trivial
    . trivial
    . exact <h.1.symm, h.2.symm>
  loopless := by
    constructor
    rintro (i | x)
    . exact fun h => h rfl
    . exact fun h => B.loopless.irrefl _ h.2
noncomputable def closedNeighborhood {V : Type*} [Fintype V]
    (B : SimpleGraph V) (v : V) : Finset V := by
  classical
  exact insert v (B.neighborFinset v)
noncomputable def minimumSize {V : Type*} [Fintype V] (B : SimpleGraph V) : Nat := by
  classical
  exact B.minDegree + 1
noncomputable def minimumCount {V : Type*} [Fintype V] (B : SimpleGraph V) : Nat := by
  classical
  exact ((Finset.univ.image (closedNeighborhood B)).filter
    (fun S => S.card = minimumSize B)).card
noncomputable def rootRadius {V : Type*} [Fintype V] (G : SimpleGraph V) : Real :=
  sSup {r : Real | exists z : Complex, (dominationPolynomial G).eval z = 0 /\ norm z = r}
noncomputable def scale (a c s m : Nat) : Real :=
  ((c : Real) * s * m / (a * Real.log m)) ^ (1 / (c : Real))
def Attainment : Prop :=
  forall {V : Type} [Fintype V] [Nonempty V] (B : SimpleGraph V)
    (a m : Nat), 0 < a -> 0 < m ->
      IsGreatest {r : Real | exists z : Complex,
        (dominationPolynomial (petalJoin a m B)).eval z = 0 /\ norm z = r}
        (rootRadius (petalJoin a m B))
def GeneralAsymptotic : Prop :=
  forall {V : Type} [Fintype V] [Nonempty V] (B : SimpleGraph V)
    (a : Nat), 0 < a ->
      Tendsto (fun m : Nat =>
        rootRadius (petalJoin a m B) / scale a (minimumSize B) (minimumCount B) m)
        atTop (nhds 1)
abbrev friendshipGraph (m : Nat) := petalJoin 1 m (top : SimpleGraph (Fin 2))
def FriendshipAsymptotic : Prop :=
  Tendsto (fun m : Nat => rootRadius (friendshipGraph m) /
    Real.sqrt (2 * (m : Real) / Real.log m)) atTop (nhds 1)
\end{verbatim}}
\noindent Here \texttt{Complex[X]} is the ring of complex polynomials,
\texttt{sum} a sum over all subsets of the vertex set, \texttt{(+)} and
\texttt{x} stand for the disjoint sum $\oplus$ and the product $\times$ of
types, \texttt{.} for the bullet $\cdot$, \texttt{<\,>} for the anonymous
constructor $\langle\,\rangle$, \texttt{!=} for $\ne$, \texttt{norm z} for
$\|z\|$, \texttt{nhds 1} for Mathlib's $\mathcal N\,1$, the neighbourhood filter
of~$1$, and \texttt{top} for
$\top$, the complete graph. \texttt{Dominates G S} is the definition of a
dominating set, and \texttt{dominationPolynomial G} is $D(G,x)$, the sum of
$x^{|S|}$ over the dominating sets $S$ of the actual graph. The vertex set of
\texttt{petalJoin a m B} is the disjoint sum of $\{0,\dots,a-1\}$ (the centre)
and the pairs (copy index, vertex of $B$); two centre vertices are adjacent
when distinct, centre and petal vertices are always adjacent, and two petal
vertices are adjacent when they lie in the same copy and are adjacent in $B$.
So \texttt{petalJoin a m B} is $J_m(a,B)$, and \texttt{friendshipGraph m} is
$K_1\vee mK_2=F_m$. \texttt{minimumSize B} is $c=1+\delta(B)$ (the vertex set
is nonempty, so the minimum degree is the true minimum), and
\texttt{minimumCount B} first takes the image of $v\mapsto N_B[v]$, which
removes repetitions, and then counts the sets of size $c$, which is $s$.
\texttt{scale a c s m} is $S_m$, with the natural logarithm and the real
power; its values for $m\le1$ do not affect a limit as $m\to\infty$.
\texttt{rootRadius G} is the supremum of the moduli of all complex roots, not
only of the real ones. \texttt{Attainment} is Theorem~\ref{thm:attain},
\texttt{GeneralAsymptotic} is Theorem~\ref{thm:main} and
\texttt{FriendshipAsymptotic} is Corollary~\ref{cor:friend}, stated directly
with $\sqrt{2m/\log m}$.

The final reviewer (see the last section) checked three points.
(1)~\emph{The supremum is a maximum.} On its own, a supremum of a set of
reals could take a default value; \texttt{Attainment} asserts, for every
$a,m\ge1$ and every $B$, that \texttt{rootRadius} is the greatest element of
the set of moduli of roots, so in \texttt{GeneralAsymptotic} it is the
modulus of an actual root. Together the two targets give both directions of
the $\varepsilon$-form of Theorem~\ref{thm:main}, including the existence of
a root of modulus at least $(1-\varepsilon)S_m$; the reviewer derived this
unfolded form for friendship graphs from the frozen targets in its own audit
file. (2)~\emph{The constant.} \texttt{minimumCount} counts distinct closed
neighbourhoods, not vertices; for $B=K_2$ or $K_3$ the two counts differ, and
the reviewer's numerical checks separate them. (3)~\emph{Quantifiers.}
\texttt{GeneralAsymptotic} holds for every finite nonempty vertex type, every
simple graph $B$ on it (disconnected graphs and isolated vertices allowed)
and every $a\ge1$, with no further hypotheses.

\paragraph{Review and freezing.}
\texttt{Challenge.lean} was reviewed against~\cite{abj} by a Codex agent that
had not written it, and frozen, by recording its SHA-256 hash, on 20 November
2026 at 18:06~UTC. The development logs record no compilation of a proof file
before that time, and the file was not changed afterwards.

\paragraph{The proofs.}
In \texttt{Proof.lean}, \texttt{Proof.Attainment} is \texttt{attainment};
\texttt{Proof.GeneralAsymptotic} combines \texttt{lower\_\allowbreak bound}
(file \texttt{LowerBound.lean}) and \texttt{upper\_\allowbreak bound} (file
\texttt{Proof.lean}, from \texttt{roots\_\allowbreak upper} and
\texttt{root\_\allowbreak equation}) with Mathlib's
\texttt{tendsto\_\allowbreak order}; and \texttt{Proof.FriendshipAsymptotic}
is \texttt{friendship\_\allowbreak of\_\allowbreak general} applied to it,
which computes $c=2$ and $s=1$ for $K_2$. Table~\ref{tab:lean} lists the Lean
counterparts of the steps of the written proof (names in the namespace
\texttt{TaskP.GraphRootRadius.Proof}).

\begin{table}[ht]
\centering\footnotesize
\renewcommand{\arraystretch}{1.0}
\begin{tabular}{@{}>{\raggedright\arraybackslash}p{0.40\textwidth}>{\raggedright\arraybackslash}p{0.56\textwidth}@{}}
\hline
Step of the written proof & Lean (file) \\
\hline
Theorem~\ref{thm:attain} &
\texttt{top\_coefficient}, \texttt{poly\_ne\_zero}, \texttt{zero\_root},
\texttt{attain\_graph}, \texttt{attainment} (\texttt{RootData}) \\
Lemma~\ref{lem:join}: dominating sets of $J_m$ &
\texttt{dominates\_join\_iff}, \texttt{dominates\_nonempty}, with
\texttt{bundle}, \texttt{familyEquiv}, \texttt{card\_bundle}
(\texttt{JoinDecomposition}) \\
Lemma~\ref{lem:join}: the formula & \texttt{eval\_join}
(\texttt{JoinPolynomial}) \\
Bad sets and complements & \texttt{Misses}, \texttt{misses\_iff},
\texttt{bad}, \texttt{min\_le\_misses}, \texttt{bad\_count}
(\texttt{PetalCoefficient}) \\
Lemma~\ref{lem:Q}(a), (b) & \texttt{normalized} (\texttt{PetalCoefficient});
\texttt{normalized\_eval}, \texttt{normalized\_at\_inverse}
(\texttt{PetalEvaluation}) \\
Lemma~\ref{lem:Q}(c), and $s\ge1$ & \texttt{normalized\_tail}
(\texttt{PetalCoefficient}); \texttt{minimumCount\_pos}
(\texttt{PetalEvaluation}); in the form $Q=1+\Delta$:
\texttt{normalized\_eq\_one\_delta} (\texttt{LowerBound}), with
\texttt{delta} (\texttt{LogLimit}) \\
Lemma~\ref{lem:rooteq} & \texttt{root\_equation} (\texttt{JoinPolynomial}) \\
Lemma~\ref{lem:scale} & \texttt{scale\_pos}, \texttt{scale\_pow},
\texttt{scale\_log\_limit}, \texttt{scale\_atTop} (\texttt{Scale}) \\
Proposition~\ref{prop:upper} & \texttt{upper\_Q},
\texttt{upper\_separation}, \texttt{roots\_upper} (\texttt{UpperEstimate});
\texttt{upper\_bound} (\texttt{Proof}) \\
The functions $H_m$ and Lemma~\ref{lem:limit} & \texttt{warg},
\texttt{residual}, \texttt{logSlope}, \texttt{log\_id},
\texttt{warg\_limit}, \texttt{delta\_limit}, \texttt{scaled\_log},
\texttt{residual\_limit} (\texttt{LogLimit}) \\
Lemma~\ref{lem:dip} & \texttt{zero\_of\_dip}, \texttt{compact\_eventually},
\texttt{eventual\_zero\_of\_dip} (\texttt{DiskZero}) \\
Choice of the disc; holomorphy; zeros of $H_m$ & \texttt{root\_disk},
\texttt{good\_disk}, \texttt{residual\_diff}, \texttt{residual\_zeros}
(\texttt{RootDisk}) \\
Lemma~\ref{lem:recover}; Proposition~\ref{prop:lower} &
\texttt{recover\_root}, \texttt{lower\_bound} (\texttt{LowerBound}) \\
Theorem~\ref{thm:main}; Corollary~\ref{cor:friend} &
\texttt{Proof.GeneralAsymptotic}, \texttt{Proof.FriendshipAsymptotic}
(\texttt{Proof}); \texttt{friendship\_of\_general} (\texttt{RootData}) \\
\hline
\end{tabular}
\caption{Lean counterparts of the steps of the proof.}
\label{tab:lean}
\end{table}

\noindent In Lean, $H_m$ is \texttt{residual} and $w_m(v)$ is
\texttt{warg}; the principal logarithm, the finiteness of the root set of a
nonzero polynomial, the existence of $\omega$ (algebraic closedness of $\C$)
and the maximum modulus principle are taken from Mathlib. The test file
\texttt{CheckRootRadius.lean} restates the three targets with the frozen
types (\texttt{check\_\allowbreak Attainment},
\texttt{check\_\allowbreak GeneralAsymptotic},
\texttt{check\_\allowbreak FriendshipAsymptotic}). The axioms of these three
theorems and of \texttt{Proof.Attainment}, \texttt{Proof.GeneralAsymptotic}
and \texttt{Proof.FriendshipAsymptotic}, printed by
\texttt{CheckRootRadius.lean}, are exactly \texttt{propext},
\texttt{Classical.choice} and \texttt{Quot.sound}. No file uses
\texttt{sorry}, \texttt{native\_\allowbreak decide},
\texttt{bv\_\allowbreak decide} or an \texttt{axiom} declaration; the one use
of \texttt{decide} (proving $0<1$ for natural numbers) is checked by the
kernel.

\paragraph{Checks.}
The $14$ files ($11$ proof modules, \texttt{Challenge.lean},
\texttt{Proof.lean} and \texttt{CheckRootRadius.lean}) were compiled module
by module in two new empty directories on the author's Mac, with the pinned
toolchain and with warnings treated as errors, on 20 November 2026
(18:24--18:30~UTC), and every module was replayed with \texttt{leanchecker};
each directory also ran one build of the whole package. All $58$ commands
($29$ per directory) ended with exit code~$0$, the \texttt{leanchecker} runs
printed nothing, and the two directories produced identical \texttt{.olean}
files. The final reviewer (see the next section) compiled the $14$ files
again, together with a short audit file of its own, in a third new directory
on the author's Mac on 20 November 2026 (about 20:53--20:56~UTC, according to
the times of its output files), with its own compilation script and without
the project's build tool, and replayed every module with
\texttt{leanchecker}. Every reviewed module compiled and replayed with exit
code~$0$, with no errors or warnings reported; a first draft of the audit
file contained an error in the reviewer's own example and was corrected and
recompiled, and no reviewed file was changed. Mathlib and its dependencies
were taken from a shared prebuilt cache at the pinned commits and were not
rechecked; the \texttt{leanchecker} replays cover the $14$ project modules.
In the repository the files have module names under
\texttt{Research.TaskP.RootRadius} and differ from the checked files only in
their \texttt{import} lines (\texttt{Challenge.lean} and
\texttt{DiskZero.lean}, which import only Mathlib, are unchanged).
The repository versions were made by a script that changed only the local \texttt{import} lines, and on 20 November 2026 (21:04--21:16~UTC) they were compiled and replayed with \texttt{leanchecker} again, module by module, in two new empty directories on the author's Mac, one after the other, together with the files of two other papers of version~\repoversion{} ($48$ modules in all); there were no warnings, and the two directories produced identical \texttt{.olean} files and the same axiom lists. In the whole-project rebuild of version~\repoversion{} of the repository (on the author's Windows computer, in two passes on 20 and 21 November 2026), these files were compiled and replayed with \texttt{leanchecker} again in pass~2 (21 November 2026), which built the files not included in pass~1 on top of its results; no module needed a second attempt (details in the repository). The formalization has not been checked with
comparator or with an independent kernel.

\par}

\section*{How this paper was produced}
This work was carried out from 20 to 22 November 2026 (UTC) in a workflow directed by
the author, using OpenAI Codex with the model GPT-6.2-Astra (reasoning level
ultra) and Anthropic Claude Code with the model Claude Opus~5.5. The author
asked for a further search for substantial open problems, and Claude Code
passed this request to Codex at 17:23~UTC. Between 17:23 and 17:53~UTC, Codex
recorded $111$ problems from the literature and wrote proofs for five
candidates. For the problem of this paper, the formula for the leading
missing term (Lemma~\ref{lem:Q}(c)) was checked by exact computation on all
$1{,}099$ labelled graphs with $1$ to $5$ vertices, the join formula on $300$
small instances, and, for $F_8$, $F_{32}$ and $F_{128}$, the presence of a
non-real root in an explicit small disc was certified in integer
arithmetic; the proof was read by another Codex agent (same model)
that had not taken part in it. The author had asked in advance that the five
most important candidates be formalized at once, without waiting for a
choice; this problem was among the five selected at 18:00~UTC. From then on,
Codex agents (same model) checked the reading of the sources, searched again
for prior work, wrote the statement file, the Lean proof and the written
proof, and ran the checks. The reading of the sources was checked, and the
statement file reviewed before it was frozen, by one Codex agent that had not
taken part in the proof or in drafting the statement. According to the file
timestamps and the development logs, the proof files were created after the
freeze and compiled between 18:07 and 18:24~UTC. They follow the written
proof prepared in the search for problems, except that the existence of roots
of the predicted size is obtained from the maximum modulus principle
(Lemma~\ref{lem:dip}) instead of the contraction-mapping argument of that
draft. After the checks, an independent Claude Code agent (model Claude
Opus~5.5) that had written none of the files gave the final review. It read
the arXiv version of~\cite{abj} and the relevant section of~\cite{ra} itself,
compared the statement file with them, read the written proof, the proof
files and the acceptance records, checked the formula of
Lemma~\ref{lem:join} against brute-force enumeration of dominating sets for
nine small petals, computed root moduli for several petals in Python
(including the values in Table~\ref{tab:num} for $m\le80$ and the Newton
values for $m\ge10^4$), and rebuilt the formalization as described above. Its
decision was APPROVE\_FORMAL, with six minor remarks on the written account
(two optional), all taken into account here. Claude Code (Claude Opus~5.5)
checked the sources, including the journal version of~\cite{abj}, recomputed
the values in Table~\ref{tab:num}, and wrote this text from Codex's
manuscript and the Lean files. An independent Claude Code agent (model Claude Opus~5.5) that had written none of the files reviewed this paper, and another reviewed the release; their corrections were applied. Throughout, ``we'' refers to
this workflow; the reading and searching described in the paper were done by
the agents. Codex agents searched for prior work on 20 November 2026, in the
search for problems and again before the statement was frozen; they found
the bounds of~\cite{ra}, but no sharp asymptotic for the largest domination
roots of friendship graphs and no matching lower bound (details in the
repository). Claude Code searched again on 21 November 2026 (17:54--18:21~UTC; the arXiv, Semantic Scholar, Crossref, DataCite, Zenodo, MathDB, GitHub and repositories of AI-assisted mathematics) and found no work on the largest domination roots of friendship graphs; the closest new finds, a note on Zenodo on the real domination roots of book graphs and a Lean formalization of the general bound $|z|\le n$ for the domination roots of graphs of order~$n$, do not overlap with the results here. Claude Code searched again on 22 November 2026 (05:39--05:51~UTC), when the OpenAlex lists of citing works refused on 21 November were retrieved, and found nothing new. This reports the coverage of the searches,
not a guarantee that they were complete. The author is not a professional
mathematician. The author has read this paper and, using a side-by-side table, compared its Lean statements with the main theorems: Theorems~\ref{thm:attain} and~\ref{thm:main} and Corollary~\ref{cor:friend}. No human mathematician has reviewed this paper.
The correctness of the results rests on the Lean proofs and on the agreement
of the Lean statements with the question of~\cite{abj}; the written proofs
are given for the reader. AI systems are not authors of this paper; the
author takes full responsibility for its content.

\medskip
\noindent Thanks are due to Saeid Alikhani, Jason I. Brown and Somayeh Jahari
for the question.

{\small
\begin{thebibliography}{9}
\setlength{\itemsep}{0pt}
\bibitem{abj} S.~Alikhani, J.~I. Brown and S.~Jahari, On the domination
polynomials of friendship graphs, \emph{Filomat} 30 (2016), no.~1, 169--178,
\url{https://doi.org/10.2298/FIL1601169A}; arXiv:1401.2092v2 (2014).
\bibitem{ap} S.~Alikhani and Y.~H. Peng, Introduction to domination
polynomial of a graph, \emph{Ars Combin.} 114 (2014), 257--266;
arXiv:0905.2251v1 (2009). Theorem numbers refer to the arXiv version.
\bibitem{bhsw} J.~I. Brown, C.~Hickman, A.~D. Sokal and D.~G. Wagner, On the
chromatic roots of generalized theta graphs, \emph{J. Combin. Theory Ser. B}
83 (2001), no.~2, 272--297, \url{https://doi.org/10.1006/jctb.2001.2057};
arXiv:math/0012033.
\bibitem{repo} H.~Dong, \emph{ai4math-results}, repository,
\url{https://github.com/Horace-Maxwell/ai4math-results}, version~\repoversion{}
(2026); archived at Zenodo under the concept DOI
\url{https://doi.org/10.5281/zenodo.22970254}.
\bibitem{kpstt} T.~Kotek, J.~Preen, F.~Simon, P.~Tittmann and M.~Trinks,
Recurrence relations and splitting formulas for the domination polynomial,
\emph{Electron. J. Combin.} 19 (2012), no.~3, \#P47; arXiv:1206.5926v2.
\bibitem{mathlib}
The mathlib Community, The Lean mathematical library, in \emph{Proc.\ 9th ACM
SIGPLAN Int.\ Conf.\ Certified Programs and Proofs} (CPP '20), ACM, 2020,
\url{https://doi.org/10.1145/3372885.3373824}.
\bibitem{ra} B.~A. Rather, On roots of domination polynomials for friendship
and book graphs, arXiv:2604.08998v2, 2026.
\end{thebibliography}}
\end{document}
